\chapter{异常控制流}

异常控制流（exceptional control flow，ECF）是对处理器状态变化的响应：
控制不再按下一条指令前进，而是转入操作系统内核或另一个进程。
它实现部分在硬件、部分在操作系统。
本章对应课堂讲义 \texttt{14-ecf-procs}（异常、进程与进程控制）；
信号、非本地跳转与相关工具尚未讲授，对应各节仍只保留目录。

\section{异常}

异常（exception）是对事件（event，处理器状态的变化）的响应：
控制从应用程序突然转到操作系统内核（kernel）。
内核是操作系统中常驻内存的那一部分，不是 \code{top} 里与用户程序并列的一个进程。
事件可以是当前指令直接引起的（除零、算术溢出、缺页），
也可以与当前指令无关（定时器到期、I/O 完成、键入 Ctrl-C）。

% BOOKFIG-BEGIN fig-8.1-exception.png
\bookfig{fig-8.1-exception.png}{8.1}{异常的剖析：事件把控制从应用程序转到异常处理程序。}{$I_{\mathrm{current}}$：发生事件时正在执行的指令；$I_{\mathrm{next}}$：若无异常则将执行的下一条；handler：内核中处理该事件的子程序}
% BOOKFIG-END fig-8.1-exception.png

处理程序结束后三条路：
回到 $I_{\mathrm{current}}$、回到 $I_{\mathrm{next}}$，或中止被中断的程序。
输入是当前指令与事件，输出是这三种控制转移之一。

\subsection{异常处理}

系统启动时，操作系统分配并初始化一张跳转表，称为异常表（exception table），
又称中断向量。
表项 $k$ 存放异常号为 $k$ 的处理程序入口地址。
运行时处理器检测到事件，确定异常号 $k$，经异常表做间接调用。

% BOOKFIG-BEGIN fig-8.2-extable.png
\bookfig{fig-8.2-extable.png}{8.2}{异常表：第 $k$ 项是异常 $k$ 的处理程序地址。}{$k$：异常号，非负整数下标；handler $k$：对应内核子程序的入口}
% BOOKFIG-END fig-8.2-extable.png

\[
\operatorname{Handler}(k)=\operatorname{Table}[k],
\qquad
k\in\{0,1,\ldots,n-1\}.
\]
输入 $k$ 由硬件提供：同步异常由 CPU 根据指令译码与执行结果生成，
异步中断由中断控制器放到总线上。
输出 $\operatorname{Handler}(k)$ 是内核代码的虚拟地址，不是用户程序里的函数指针。
$k$ 的取值范围是该体系结构异常表的合法下标，例如 x86 上缺页、除零各有固定编号。

\subsection{异常的类别}

异常分为四类。图~8.4 的属性写成
\begin{center}
\begin{tabular}{llll}
\toprule
类别 & 起因 & 异步/同步 & 返回行为 \\
\midrule
中断 & I/O 设备信号 & 异步 & 总是回到 $I_{\mathrm{next}}$ \\
陷阱 & 有意的异常 & 同步 & 总是回到 $I_{\mathrm{next}}$ \\
故障 & 可能恢复的错误 & 同步 & 可能回到 $I_{\mathrm{current}}$ \\
终止 & 不可恢复的错误 & 同步 & 永不返回 \\
\bottomrule
\end{tabular}
\end{center}
\noindent{\heiti 图~8.4}\quad 异常的类别。\\
符号：同步=由正在执行的指令直接引起；异步=来自处理器外部。

中断（interrupt）由外部设备拉高中断脚并提供异常号。
定时器中断使内核能从用户程序收回 CPU，用于调度。
处理完回到 $I_{\mathrm{next}}$。

陷阱（trap）是有意执行的指令，例如系统调用与断点。
处理完回到 $I_{\mathrm{next}}$，对用户像一次进入内核的过程调用。

故障（fault）由当前指令无意触发，但内核或许能补救。
缺页常把那一页调入物理内存后重做 $I_{\mathrm{current}}$；
保护错往往不可恢复，处理程序中止进程。
浮点异常也属此类。

终止（abort）由不可恢复错误引起，例如硬件校验错。
处理程序中止当前程序，不回到用户代码。

\subsection{Linux/x86-64系统中的异常}

x86-64 Linux 上，故障与陷阱的例子包括：除零、缺页、通用保护、
以及 \code{syscall} 指令引起的系统调用。
系统调用是陷阱：用户程序把编号放入约定寄存器再执行 \code{syscall}。

\begin{center}
\begin{tabular}{cll}
\toprule
编号 & 名字 & 作用 \\
\midrule
0 & \code{read} & 读文件 \\
1 & \code{write} & 写文件 \\
2 & \code{open} & 打开文件 \\
3 & \code{close} & 关闭文件 \\
4 & \code{stat} & 取文件信息 \\
57 & \code{fork} & 创建进程 \\
59 & \code{execve} & 执行程序 \\
60 & \code{\_exit} & 结束进程 \\
62 & \code{kill} & 向进程发信号 \\
\bottomrule
\end{tabular}
\end{center}
\noindent{\heiti 图~8.10}\quad Linux x86-64 上若干系统调用。\\
符号：编号=进入内核时 \texttt{\%rax} 的值；名字=C 库包装函数。

编号是内核约定，与缺页等异常号不是同一张日常清单。
C 中写 \code{open}、\code{fork}，最终落到这些编号之一。

% BOOKFIG-BEGIN fig-8.6-syscall.png
\bookfig{fig-8.6-syscall.png}{8.6}{系统调用示例：\texttt{open} 经 \texttt{syscall} 进内核，返回后继续下一条。}{\texttt{\%rax}：进去时为调用号 $2$，回来时为返回值；$I_{\mathrm{next}}$：\texttt{syscall} 之后的比较指令}
% BOOKFIG-END fig-8.6-syscall.png

x86-64 Linux 约定：进入时 \code{\%rax} 为调用号，
\code{\%rdi}、\code{\%rsi}、\code{\%rdx}、\code{\%r10}、\code{\%r8}、\code{\%r9} 为参数；
返回后 \code{\%rax} 为结果，负值表示出错。
输入是调用号与参数寄存器，输出是返回值标量。
该标量对 \code{open} 是文件描述符或错误码的相反数一类约定。

% BOOKFIG-BEGIN fig-8.7-fault.png
\bookfig{fig-8.7-fault.png}{8.7}{故障示例：合法地址缺页时拷页并重做 \texttt{movl}。}{$I_{\mathrm{current}}$：引起缺页的写指令；disk$\to$mem：内核把该页调入物理内存}
% BOOKFIG-END fig-8.7-fault.png

对 \code{a[500]=13}，虚拟地址合法但页不在内存：同步故障。
内核拷页后回到 $I_{\mathrm{current}}$ 重做同一条 \code{movl}。
对 \code{a[5000]=13}，硬件仍先报缺页，内核发现地址无合法映射，
向该进程发 \code{SIGSEGV}，默认终止，不再重做。
$\operatorname{SIGSEGV}$ 是段错误对应的信号编号（Linux 上常为 $11$），
表示当前进程访问了不该访问的地址或权限不够。

\section{进程}

进程（process）是一个正在运行的程序的实例（instance）。
它不是磁盘上的可执行文件，也不是处理器硬件。
同一份程序可以同时对应多个进程；一个核可以轮流跑多个进程。

% BOOKFIG-BEGIN fig-8.13-aspace.png
\bookfig{fig-8.13-aspace.png}{8.13}{进程提供逻辑控制流与私有地址空间两层抽象。}{CPU/Registers：该进程眼里的处理器；Stack/Heap/Data/Code：该进程的虚拟存储器分区}
% BOOKFIG-END fig-8.13-aspace.png

进程给每个程序两层关键抽象：逻辑控制流（好像独占 CPU）与私有地址空间（好像独占主存）。
前者由上下文切换提供，后者由虚拟内存提供。

\subsection{逻辑控制流}

处理器从开机到关机执行的指令序列是物理控制流。
每个进程对应其中一段交错后的投影，称为逻辑控制流：
进程自己看见的指令一条接一条，中间插入的别的进程被抽象掉。
上下文切换把当前进程的寄存器等保存到内存，装入另一个进程的现场，
使每个程序以为 CPU 只在执行自己。

\subsection{并发流}

两个进程并发（concurrent），当且仅当它们的逻辑流在时间上重叠；
否则称为顺序（sequential）。
并发不要求同一时刻两条用户指令都在占用同一颗核。

% BOOKFIG-BEGIN fig-8.12-concurrent.png
\bookfig{fig-8.12-concurrent.png}{8.12}{单核上三个进程的交错：流的生存期重叠则并发。}{Time：向下为时间；竖杠：该时间片内核在执行哪个进程}
% BOOKFIG-END fig-8.12-concurrent.png

设进程 $A,B,C$ 的生存区间为 $[\alpha_X,\omega_X]$。
\[
\operatorname{Concurrent}(X,Y)
\iff
[\alpha_X,\omega_X]\cap[\alpha_Y,\omega_Y]\neq\emptyset.
\]
输入是两条流的起止时刻，输出是真假。
真表示二者并发，假表示顺序。
单核示例中 $A$ 与 $B$、$A$ 与 $C$ 重叠故并发，$B$ 在 $C$ 开始前结束故顺序。

物理上，单核并发流的执行片段在时间上不相交；
用户视角仍可把各流想成彼此并行。
多核时每个核可同时执行一个不同进程，调度仍由内核完成。

\subsection{私有地址空间}

每个进程有自己的虚拟地址空间：代码、数据、堆、栈。
相同的虚拟地址在不同进程中映射到不同的页。
这解释了缺页非法时只终止当前 PID，以及 \code{fork} 之后父子对 \code{x} 的修改互不影响。

\subsection{用户模式和内核模式}

处理器用模式位区分用户模式与内核模式。
用户模式不能执行特权指令、不能直接改页表或访问设备。
异常把控制转入内核并把模式改为内核模式；
处理程序作为某个已存在进程的一部分运行，内核不是单独的进程。
从内核返回用户代码时恢复用户模式。

\subsection{上下文切换}

% BOOKFIG-BEGIN fig-8.14-cswitch.png
\bookfig{fig-8.14-cswitch.png}{8.14}{控制流在内核中从进程 $A$ 转到进程 $B$。}{user code：用户指令；kernel code：内核中保存/装入寄存器并切换地址空间的那段}
% BOOKFIG-END fig-8.14-cswitch.png

一次上下文切换包含：把当前寄存器写入该进程的保存区，
调度下一个进程，装入其保存的寄存器并切换地址空间（换页表）。
切换发生在内核中，由系统调用、缺页或定时器中断等异常触发。
单核上任意时刻只有一套活寄存器；未执行进程的寄存器值存在内存里。

\section{系统调用错误处理}

Linux 系统级函数失败时通常返回 $-1$，并设置全局变量 \code{errno} 表示原因。
必须检查每个非 \code{void} 系统级函数的返回值。
\[
\operatorname{Fail}(f)\iff r=-1,
\]
此时 \code{strerror}(\code{errno}) 把编号译成字符串。
$r$ 是函数返回的标量；$-1$ 表示失败，非负（或约定的成功值）表示成功。

课堂把打印与退出收成
\begin{lstlisting}
void unix_error(char *msg)
{
    fprintf(stderr, "%s: %s\n", msg, strerror(errno));
    exit(0);
}
\end{lstlisting}
再把调用与检查包成 Stevens 风格包装：大写 \code{Fork} 内调 \code{fork}，
失败则 \code{unix\_error}，成功才把 \code{pid} 返回给调用者。
小写名为系统函数，大写名为课件包装，不是内核另提供的系统调用。

\section{进程控制}

程序员通过系统调用创建、终止、回收进程并加载新程序。

\subsection{获取进程ID}

\[
p=\operatorname{getpid}(),
\qquad
p_{\mathrm{par}}=\operatorname{getppid}().
\]
输入为空，输出是 \code{pid\_t} 整数。
$p$ 是当前进程的 PID，$p_{\mathrm{par}}$ 是父进程的 PID，取值在内核分配的正整数编号中。
PID $0$ 不是用户进程编号。

\subsection{创建和终止进程}

从程序员视角，进程处于三态之一。
Running：正在执行，或等待被调度且终将被内核选中（含等待 I/O）。
Stopped：挂起，在收到继续通知前不会被调度。
Terminated：永久停止。

进程因下列原因之一变为 Terminated：
收到默认动作为终止的信号（如 \code{SIGSEGV}）、从 \code{main} 返回、或调用 \code{exit}。
\[
\operatorname{exit}(\textit{status})
\]
结束当前进程并留下退出码 \textit{status}。
约定 $0$ 为正常，非 $0$ 为出错。
从 \code{main} 返回整数 $n$ 等价于 \code{exit}($n$)。
\code{exit} 调用一次、永不返回。

父进程调用 \code{fork} 创建新的 Running 子进程。
\[
r=\operatorname{fork}().
\]
输入为空。输出 $r$ 在两个进程中各出现一次：
父进程中 $r$ 等于孩子的 PID（正整数），子进程中 $r=0$，失败时 $r=-1$。
$0$ 只出现在孩子对返回值的使用中，不是孩子的 PID；
孩子的真实编号由 \code{getpid} 给出。

孩子得到父进程虚拟地址空间的一份相同但分离的拷贝，以及打开文件描述符的拷贝，
并具有不同的 PID。
\code{fork} 调用一次、返回两次。

\begin{lstlisting}
int x = 1;
pid = Fork();
if (pid == 0) {          /* Child */
    printf("child : x=%d\n", ++x);
    exit(0);
}
printf("parent: x=%d\n", --x);
exit(0);
\end{lstlisting}
\code{Fork} 返回时两份 $x$ 都是 $1$。
孩子执行 ${++}x$ 印 $2$，父亲执行 ${--}x$ 印 $0$，之后各自独立。
谁先 \code{printf} 不能事先保证。
孩子必须 \code{exit}，否则会落到父亲的打印语句。
两边的 \code{stdout} 常仍指向同一终端。

进程图（process graph）刻画并发语句的偏序：
顶点是一次语句执行，边 $a\to b$ 表示 $a$ 一定发生在 $b$ 之前。
任意拓扑序对应一种可行的总顺序。

% BOOKFIG-BEGIN fig-8.16-pgraph.png
\bookfig{fig-8.16-pgraph.png}{8.16}{\texttt{fork.c} 的进程图：\texttt{fork} 处分叉，两支 \texttt{printf} 之间无箭头。}{$x{=}1$：分叉前的变量；child $x{=}2$ / parent $x{=}0$：两支上的输出}
% BOOKFIG-END fig-8.16-pgraph.png

连续两次无 \code{if} 的 \code{fork} 使每个进程都继续并再分叉：
\code{L0} 印 $1$ 次，\code{L1} 印 $2$ 次，\code{Bye} 印 $4$ 次。
每个 \code{Bye} 顶点都在该进程自己的 \code{L1} 之后，故全局第一个 \code{Bye} 前至少有一个 \code{L1}。

\code{if (fork() != 0)} 只让父进程继续嵌套，得到 $3$ 个进程，
\code{L0}、\code{L1}、\code{L2} 各 $1$ 次，\code{Bye} 共 $3$ 次；
最初父进程的路径含 \code{L2}$\to$\code{Bye}。
\code{if (fork() == 0)} 只让子进程继续嵌套，次数相同，链在孩子一侧；
孙子的路径含 \code{L2}$\to$\code{Bye}。

\subsection{回收子进程}

进程终止后仍占用退出码与内核表项，称为僵尸（zombie）。
父进程用 \code{wait} 或 \code{waitpid} 回收（reap）：
\[
w=\operatorname{wait}(\&\textit{status}).
\]
输入是指向整数的指针（或 \code{NULL}）。
输出 $w$ 是被回收孩子的 PID。
若指针非空，则 \textit{status} 被写成编码后的结束原因。
用 \code{<sys/wait.h>} 中的宏解码：
\code{WIFEXITED} 为真时 \code{WEXITSTATUS} 给出 \code{exit} 的整数；
\code{WIFSIGNALED} 为真时 \code{WTERMSIG} 给出信号编号（如 \code{SIGSEGV}）。
这些宏的输入是 \textit{status} 那一个整数，输出是真假或非负编号。

\code{wait} 挂起当前进程直到它的某一个孩子终止；多个已结束的孩子按任意顺序回收。
一次调用回收一个孩子，成功 \code{fork} 了 $N$ 个就要回收 $N$ 次。

\[
w=\operatorname{waitpid}(\textit{pid},\&\textit{status},\textit{opt}).
\]
$\textit{pid}>0$ 只等该进程；$\textit{pid}=-1$ 与 \code{wait} 等价。
$\textit{opt}=0$ 表示死等结束；
\code{WNOHANG} 使无状态可报时立即返回 $0$；
\code{WUNTRACED}、\code{WCONTINUED} 分别在停止与继续时也可返回。
$w$ 为孩子 PID、在 \code{WNOHANG} 下可为 $0$，出错为 $-1$。

父进程不回收且自己也终止时，孤儿由 \code{init}（PID $1$）代收。
长命进程（shell、服务器）必须显式回收，否则僵尸占满进程表。
孩子已 \code{exit}、父亲死循环不 \code{wait} 时，\code{ps} 显示 \code{<defunct>}；
杀死父亲后 init 收掉该行。
父亲先 \code{exit}、孩子死循环时，孩子是仍在运行的孤儿，不是僵尸，须另外结束该 PID。

\code{wait} 在进程图上增加孩子 \code{exit} 到父亲 \code{wait} 的边，
因此 \code{CT}（child terminated）不能出现在孩子的打印之前。

\subsection{让进程休眠}

\pending

\subsection{加载并运行程序}

\code{execve} 在当前进程中加载并运行一个新程序：
\[
r=\operatorname{execve}(\textit{filename},\textit{argv},\textit{envp}).
\]
输入 \textit{filename} 是可执行文件或以 \texttt{\#!} 开头的脚本路径，
\textit{argv}、\textit{envp} 是以 \code{NULL} 结尾的字符串指针数组，
约定 $\textit{argv}[0]$ 为程序名。
成功时不返回（$r$ 不出现）；失败返回 $-1$。
它覆盖代码、数据与栈，保留 PID、打开的文件与部分信号上下文。
旧用户地址空间被取消映射，内核在高地址重新铺用户栈且仍向低地址增长。

% BOOKFIG-BEGIN fig-8.22-stack.png
\bookfig{fig-8.22-stack.png}{8.22}{新程序启动时用户栈的组织。}{\texttt{argc} 在 \texttt{\%rdi}；\texttt{argv} 在 \texttt{\%rsi}；\texttt{envp} 在 \texttt{\%rdx}；\texttt{environ}：指向环境表的全局指针}
% BOOKFIG-END fig-8.22-stack.png

高地址存放环境串与参数串，其下为 \code{envp[]}、\code{argv[]}，
再下为 \code{libc\_start\_main} 与将来 \code{main} 的栈帧。
x86-64 把 $\textit{argc}$、$\textit{argv}$、$\textit{envp}$ 放入
\code{\%rdi}、\code{\%rsi}、\code{\%rdx} 后调用 \code{main}。
$\textit{argc}$ 是参数个数（不含结尾 \code{NULL}），为非负整数。

\subsection{利用fork和execve运行程序}

% BOOKFIG-BEGIN fig-8.20-argv.png
\bookfig{fig-8.20-argv.png}{8.20}{孩子中 \texttt{execve("/bin/ls", myargv, environ)} 的参数与环境表。}{\texttt{myargv[0..2]}：\texttt{/bin/ls}、\texttt{-lt}、\texttt{/usr/include}；\texttt{argc}=3}
% BOOKFIG-END fig-8.20-argv.png

典型用法：父亲 \code{fork}，孩子调用 \code{execve}，父亲 \code{wait}。
\begin{lstlisting}
if ((pid = Fork()) == 0) {
    if (execve(myargv[0], myargv, environ) < 0) {
        printf("%s: Command not found.\n",
               myargv[0]);
        exit(1);
    }
}
\end{lstlisting}
成功则该 PID 的映像换成 \code{ls}，旧代码不回到 \code{printf}。
失败才印错误并以 $1$ 退出。
\code{ls} 是已经链接好的可执行文件，不必与当前程序静态链在一起。
父亲仍运行原来的程序并回收孩子；没有「把旧程序再装回该 PID」这一步。

\section{信号}

\pending

\subsection{信号术语}

\pending

\subsection{发送信号}

\pending

\subsection{接收信号}

\pending

\subsection{阻塞和解除阻塞信号}

\pending

\subsection{编写信号处理程序}

\pending

\subsection{同步流以避免讨厌的并发错误}

\pending

\subsection{显式地等待信号}

\pending

\section{非本地跳转}

\pending

\section{操作进程的工具}

\pending

\section{小结}

异常是需要非标准控制流的事件，由外部中断或内部陷阱、故障（及终止）产生。
进程在系统中同时存在多个；单核同一时刻只执行一个，
但每个进程都表现为独占处理器与私有地址空间。
创建进程调用 \code{fork}：一次调用、两次返回。
结束进程调用 \code{exit}：一次调用、不返回。
回收与等待调用 \code{wait} 或 \code{waitpid}。
加载并运行新程序调用 \code{execve}：一次调用，成功则不返回。
